\documentclass[../../main.tex]{subfiles}
\begin{document}

\tikzstyle{Node_Default}=[circle, draw, thick]
\tikzstyle{Edge_Default}=[->, thick]
\begin{problem}
    求下面两个图里面每对顶点之间的最短路径:
\end{problem}
\begin{center}
\begin{tikzpicture}
\node[Node_Default] (A1) at ( 0 , 2 ) {1};
\node[Node_Default] (B1) at ( 2 , 2 ) {2};
\node[Node_Default] (C1) at ( 4 , 2 ) {3};

\node[Node_Default] (D1) at ( 0 , 0 ) {4};
\node[Node_Default] (E1) at ( 2 , 0 ) {5};
\node[Node_Default] (F1) at ( 4 , 0 ) {6};

\node[Node_Default] (A2) at ( 6  , 2 ) {1};
\node[Node_Default] (B2) at ( 8  , 2 ) {2};
\node[Node_Default] (C2) at ( 10 , 2 ) {3};

\node[Node_Default] (D2) at ( 6 , 0 ) {4};
\node[Node_Default] (E2) at ( 8 , 0 ) {5};
\node[Node_Default] (F2) at ( 10, 0 ) {6};

\draw[Edge_Default] (A1) -- (E1) node[pos=0.8, above] {1};
\draw[Edge_Default] (B1) -- (A1) node[pos=0.5, above] {1};
\draw[Edge_Default] (B1) -- (D1) node[pos=0.2, above] {2};
\draw[Edge_Default] (C1) -- (B1) node[pos=0.5, above] {2};
\draw[Edge_Default] (C1.280) -- (F1.80) node[pos=0.5, right] {8};
\draw[Edge_Default] (D1) -- (A1) node[pos=0.5, left] {4};
\draw[Edge_Default] (D1) -- (E1) node[pos=0.5, below] {3};
\draw[Edge_Default] (E1) -- (B1) node[pos=0.5, right] {7};
\draw[Edge_Default] (F1) -- (B1) node[pos=0.4, left] {5};
\draw[Edge_Default] (F1.100) -- (C1.260) node[pos=0.5, left] {10};

\draw[Edge_Default] (A2) -- (E2) node[pos=0.8, above] {-1};
\draw[Edge_Default] (B2) -- (A2) node[pos=0.5, above] {1};
\draw[Edge_Default] (B2) -- (D2) node[pos=0.2, above] {2};
\draw[Edge_Default] (C2) -- (B2) node[pos=0.5, above] {2};
\draw[Edge_Default] (C2.280) -- (F2.80) node[pos=0.5, right] {-8};
\draw[Edge_Default] (D2) -- (A2) node[pos=0.5, left] {-4};
\draw[Edge_Default] (D2) -- (E2) node[pos=0.5, below] {3};
\draw[Edge_Default] (E2) -- (B2) node[pos=0.5, right] {7};
\draw[Edge_Default] (F2) -- (B2) node[pos=0.4, left] {5};
\draw[Edge_Default] (F2.100) -- (C2.260) node[pos=0.5, left] {10};
\end{tikzpicture}
\end{center}
\begin{solution}
    \begin{itemize}
        \item [图1:]
        可以写出其邻接矩阵:
        \[
        \begin{bmatrix}
            0      & \infty & \infty & \infty & 1      & \infty \\
            1      & 0      & \infty &  2     & \infty & \infty \\
            \infty & 2      & 0      & \infty & \infty & 8      \\
            4      & \infty & \infty & 0      & 3      & \infty \\
            \infty & 7      & \infty & \infty & 0      & \infty \\
            \infty & 5      & 10     & \infty & \infty & 0
        \end{bmatrix}
        \]
        计算下一个矩阵的公式为:
        \[
        A^{k+1}_{ij} = \min \{ A^k_{ij}, A^k_{ik} + A^k_{kj} \}
        \]
        以此可算出\(A^1 - A^6\),选取对应的最短路径, 得到最终的距离矩阵:
        \[
        A = 
        \begin{bmatrix}
            0& 8& \infty& 10& 1& \infty \\
            1& 0& \infty& 2& 2& \infty \\
            3& 2& 0& 4& 4& 8 \\
            4& 10& \infty& 0& 3& \infty \\
            8& 7& \infty& 9& 0& \infty \\
            6& 5& 10& 7& 7& 0 \\
        \end{bmatrix}
        \]
        \item [图2:]
        可以写出其邻接矩阵:
        \[
        \begin{bmatrix}
            0      & \infty & \infty & \infty & -1     & \infty \\
            1      & 0      & \infty &  2     & \infty & \infty \\
            \infty & 2      & 0      & \infty & \infty & -8      \\
            -4      & \infty & \infty & 0     &  3     & \infty \\
            \infty & 7      & \infty & \infty & 0      & \infty \\
            \infty & 5      & 10     & \infty & \infty & 0
        \end{bmatrix}
        \]
        计算下一个矩阵的公式为:
        \[
        A^{k+1}_{ij} = \min \{ A^k_{ij}, A^k_{ik} + A^k_{kj} \}
        \]
        以此可算出\(A^1 - A^6\),选取对应的最短路径, 得到最终的距离矩阵:
        \[
        A =
        \begin{bmatrix}
            0& 6& \infty& 8& -1& \infty \\
            -2& 0& \infty& 2& -3& \infty \\
            -5& -3& 0& -1& -6& -8 \\
            -4& 2& \infty& 0& -5& \infty \\
            5& 7& \infty& 9& 0& \infty \\
            3& 5& 10& 7& 2& 0 \\
        \end{bmatrix}
        \]
    \end{itemize}
\end{solution}
\end{document}